\documentclass[12pt]{article}
\usepackage[utf8]{inputenc}
\usepackage[T1]{fontenc}
\usepackage{hyperref}
\title{Sawai Language Resources for K-12 Curriculum Development}
\author{Compiled by Claude Sonnet 5.5 (claude-sonnet-5-5) for Andr\'as Kornai}
\date{}
\begin{document}
\maketitle

\section{Introduction}
Sawai (ISO 639-3 szw) is an Austronesian language of the Weda area of southern Halmahera, North Maluku, Indonesia. Reported dialects are Weda, Sawai, Foya-Mafa and Messa-Dote. It is written in the Latin alphabet. This report is based on a single round of web searches.

\section{Primary Educational Resources}
No school textbook or primer was found. Writing work reportedly slowed after the war in North Maluku began in late 1999.

\section{Linguistic and Reference Materials}
A draft Sawai dictionary is online at Webonary. Its authors say it is incomplete and invite corrections: \url{https://www.webonary.work/sawai/overview/foreword/?lang=en}.

Ronald Whistler's ``Phonology of Sawai'' (in Burquest and Laidig, eds., \emph{Phonological studies in four languages of Maluku}, SIL, 1992, pp.\ 7--32) is the main academic source. I did not see a URL for the text itself. Overviews are at \url{https://en.wikipedia.org/wiki/Sawai_language} and \url{https://omniglot.com/writing/sawai.htm}. A 2021 paper on the Sawai numeral system is at \url{https://www.atlantis-press.com/proceedings/asbam-21/125973606}, and a numerals page at \url{https://lingweb.eva.mpg.de/channumerals/Sawai.htm}. No complete grammar was found.

\section{Religious and Community Resources}
None found.

\section{Digital and Audio Resources}
The Webonary dictionary is the only digital resource found. Field recordings may exist in the PARADISEC Halmahera series (\url{https://paradisec.org.au/fieldnotes/INHALM.htm}), but I did not confirm that they include Sawai.

\section{Curriculum Development Context}
The language is used by all adults but not all young people, and no school use is known. The draft dictionary and the phonology give a starting point for an orthography-based primer.
\end{document}